\section{实操视角：设计一个训练数据管线要问什么}

\begin{importantbox}{数据管线 checklist}
\begin{enumerate}
    \item \textbf{来源}：每个 source 是 live service、public dump、licensed corpus、synthetic data 还是 user data？
    \item \textbf{权限}：是否有 license？ToS 是否允许？robots.txt 怎么处理？是否含 shadow-library 风险？
    \item \textbf{抽取}：HTML/PDF/LaTeX/code/XML 如何转文本？转换器版本和失败率是什么？
    \item \textbf{过滤}：使用手工规则、语言识别、quality classifier、toxicity filter 还是 PII filter？阈值怎么定？
    \item \textbf{去重}：document-level、paragraph-level、MinHash、Bloom filter、benchmark overlap 分别怎么做？
    \item \textbf{混合}：web、code、math、academic、books、Q\&A、multilingual、synthetic data 的比例如何决定？
    \item \textbf{审计}：能否追踪某条训练样本的 provenance？能否删除某类数据？能否复现实验？
\end{enumerate}
\end{importantbox}

\begin{knowledgebox}{数据配比不是静态表格}
Pre-training 早期可能需要最大覆盖和多样性；mid-training 可能提高代码、数学、高质量学术文本或教育文本比例；post-training 则加入 instruction、preference、tool-use、RL trajectories。数据 mix 应服务训练阶段和目标能力，而不是一次性固定。
\end{knowledgebox}

\begin{warningbox}{数据工作最容易隐藏在“清洗过”三个字里}
论文常写“we cleaned and deduplicated the data”，但这句话可能包含数百个规则、多个模型过滤器、人工审计、legal review 和大量失败样本。对复现和安全来说，清洗规则和过滤器版本与模型 architecture 一样重要。
\end{warningbox}

